\documentclass{article}
\usepackage[left=1in,right=1in,top=1in,bottom=1in]{geometry}
\usepackage{amsmath, amssymb, amsthm, verbatim, enumerate, xcolor, mathpartir}

\newtheorem{task}{Task}
\newtheorem{lemma}{Lemma}
\newtheorem{thm}{Theorem}
  

%% Inference rules
%Inference rule with a title
\newcommand{\Rule}[4][]{\ensuremath{\inferrule*[right={(#2)},#1]{#3}{#4}}}
%Inference rule with no title
\newcommand{\RuleNoLabel}[3][]{\ensuremath{\inferrule*[#1]{#2}{#3}}}

%% BNF
\newcommand{\bnfdef}{\mathop{::=}}
\newcommand{\bnfalt}{\mathbin{|}}

%% XML
\newcommand{\isxml}[1]{#1~\ensuremath{\mathsf{is~XML}}}
\newcommand{\xml}[1]{\text{\texttt{#1}}}

%% RB Trees
\newcommand{\tree}{\ensuremath{T}}
\newcommand{\leaf}{\ensuremath{\mathsf{Leaf}}}
\newcommand{\bnode}[2]{\ensuremath{\mathsf{BNode}(#1, #2)}}
\newcommand{\rnode}[2]{\ensuremath{{\color{red}{\mathsf{RNode}}}(#1, #2)}}
\newcommand{\bc}{\ensuremath{\mathsf{Black}}}
\newcommand{\rc}{\ensuremath{{\color{red}{\mathsf{Red}}}}}
\newcommand{\istree}[2]{~\ensuremath{\mathsf{isTree}~#1~#2}}
\newcommand{\colr}{\ensuremath{c}}
\newcommand{\nodes}[1]{\ensuremath{\mathit{Nodes}(#1)}}


\title{CSE5095: Types and Programming Languages\\
  {\large Homework 1: Inference rules, derivations, rule induction}}
\author{Your name here}

\begin{document}
\maketitle

\section{Logistics}
\textbf{Important note: This document DOES NOT have all of the information
  you need to complete the assignment. It's just a template for typesetting
  your answers. You MUST read the writeup pdf.}

Great to see you're (at least thinking about) writing up your solution in
LaTeX though! Thanks! See HuskyCT for resources if you need help.
Some useful macros are defined in the file \texttt{macros.tex}, and
this source file gives examples of how to use them.

\section{XML}

\begin{task}\label{concat3}
Add another inference rule for showing that the concatenation of
3 XML strings (e.g., $XYZ$) is also an XML string.
\end{task}

\textbf{Answer:}

\[
\Rule{IdentityRule}
     {\isxml{X}\\
       \isxml{Y}\\
       \isxml{Z}}
     {\isxml{XYZ}}
\]

\begin{task}
  Write a derivation for the conclusion
  $\isxml{XYZ}$ assuming the premises of the rule you wrote in
  Task~\ref{concat3}.
    Make sure to include the names/numbers of the rules you're using.
\end{task}

\textbf{Answer:}

\[
\Rule{X-4}
     {\Rule{X-4}
           {\isxml{X}\\
             \isxml{Y}}
           {\isxml{XY}}\\
       \isxml{Z}}
     {\isxml{XYZ}}
\]

Typestting Hint: you can typeset a derivation by nesting more
inference rules(\texttt{\textbackslash Rule}) in the premises of an
inference rule. The code for this gets messy, so make sure you're using
indentation or other hints to keep yourself sane.
If you're already familiar with using LaTeX for math, you may be tempted to
just use \texttt{\textbackslash frac}, but this gets really difficult to
read because it makes the nested fractions smaller and smaller (and also
doesn't let you include the rule names).


\section{Red-Black Trees}
Code for the inference rules and definitions is below so you can have
examples of how to typeset these.

\[
\begin{array}{r l l l}
  Trees & \tree & \bnfdef & \leaf \bnfalt \bnode{\tree}{\tree} \bnfalt
  \rnode{\tree}{\tree}\\
  Colors & \colr & \bnfdef & \bc \bnfalt \rc\\
\end{array}
\]

{
  \centering
  \def \MathparLineskip {\lineskip=0.43cm}
  \begin{mathpar}

    \Rule{RBT-1}
         {\strut}
         {\leaf \istree{\bc}{1}}
    \and
    \Rule{RBT-2}
         {\tree_1 \istree{\colr_1}{n}\\
           \tree_2 \istree{\colr_2}{n}}
         {\bnode{\tree_1}{\tree_2} \istree{\bc}{n+1}}
    \and
    \Rule{RBT-3}
         {\tree_1 \istree{\bc}{n}\\
           \tree_2 \istree{\bc}{n}}
         {\rnode{\tree_1}{\tree_2} \istree{\rc}{n}}
  \end{mathpar}
}

\[
\begin{array}{l l l}
  \nodes{\leaf} & = & 1\\
  \nodes{\bnode{\tree_1}{\tree_2}} & = & 1 + \nodes{\tree_1} + \nodes{\tree_2}\\
  \nodes{\rnode{\tree_1}{\tree_2}} & = & 1 + \nodes{\tree_1} + \nodes{\tree_2}
\end{array}
\]

\begin{task}
  Prove by rule induction that for any tree~$\tree$,
  color~$\colr \in \{\rc, \bc\}$ and~$n \geq 1$,
  if~$\tree \istree{\colr}{n}$, then~$\nodes{\tree} \geq 2^{n-1}$.

  {\em You must} use rule induction, not any other proof technique.

  \textbf{Note:} This, together with one or two other properties of
  red-black trees, proves that a valid red-black tree is approximately
  balanced.
\end{task}

\textbf{Answer:}

Proof: By rule induction on the derivation of $\tree \istree{\colr}{n}$.
We have one case for each rule that could have concluded
$\tree \istree{\colr}{n}$.

\vspace{0.3cm}
\noindent\textbf{Case RBT-1:} $\tree = \leaf$, $\colr = \bc$, $n = 1$.

% [FILL IN: what is Nodes(Leaf)? does it satisfy Nodes(T) >= 2^(n-1)
%  when n=1? this case has no premises/IH to use.]

\vspace{0.3cm}
\noindent\textbf{Case RBT-2:} $\tree = \bnode{\tree_1}{\tree_2}$, $\colr = \bc$,
$\tree$ derived from $\tree_1 \istree{\colr_1}{m}$ and
$\tree_2 \istree{\colr_2}{m}$, with $n = m+1$.

% [FILL IN: state the induction hypothesis for T1 and T2 (careful: IH as
%  stated only applies when m >= 1 -- what happens if m = 0? is the bound
%  still true in that case?).
%  Then use Nodes(BNode(T1,T2)) = 1 + Nodes(T1) + Nodes(T2) and the IH
%  to show Nodes(T) >= 2^(n-1) = 2^m.]

\vspace{0.3cm}
\noindent\textbf{Case RBT-3:} $\tree = \rnode{\tree_1}{\tree_2}$, $\colr = \rc$,
$\tree$ derived from $\tree_1 \istree{\bc}{n}$ and
$\tree_2 \istree{\bc}{n}$.

% [FILL IN: state the IH for T1 and T2 (same n as the conclusion here,
%  unlike RBT-2). Use Nodes(RNode(T1,T2)) = 1 + Nodes(T1) + Nodes(T2)
%  and the IH to show Nodes(T) >= 2^(n-1).]

\hfill$\blacksquare$

\section{One more wrap-up question}

\begin{task}
  How long (approximately) did you spend on this homework, in total
  minutes/hours of
  actual working time?
  Your honest feedback will help me with future homeworks.
\end{task}
\textbf{Answer:}

\end{document}
